\documentclass[12pt]{article}
\usepackage{graphicx}
\usepackage{hyperref}
\usepackage{geometry}
\geometry{margin=1in}

\title{Computer Network Hardware Lab Report}
\author{Istiaq Alam}
\date{\today}

\begin{document}

\maketitle
\tableofcontents
\newpage

\section{Cable Tracer}
A \textbf{Cable Tracer} is used to trace and identify cables or wires within a group without damaging the insulation. It is useful for locating breaks or shorts in cables and identifying signal paths.

\begin{center}
%\includegraphics[width=0.5\textwidth]{https://ae01.alicdn.com/kf/HLB10vTCa6LuK1Rjy0Fhq6xpdFXaz/1-Set-Network-Line-Finder-Cable-Detector-JW-360-LAN-Network-Cable-Tester-RJ11-RJ45-UTP.jpg}
\end{center}

\section{Cable Tester}
A \textbf{Cable Tester} is a device used to test the strength and connectivity of a particular type of cable or other wired assemblies. It checks for continuity, wiring configuration, and signal loss.

\begin{center}
%\includegraphics[width=0.5\textwidth]{https://cdn.sparkfun.com//assets/parts/1/2/2/5/6/16227-CableTester__main.jpg}
\end{center}

\section{T-568B Wiring Standard}
The \textbf{T-568B} standard is a wiring scheme used for terminating network cables (Ethernet) using RJ-45 connectors. The order of colored wires is:
\begin{enumerate}
\item White Orange
\item Orange
\item White Green
\item Blue
\item White Blue
\item Green
\item White Brown
\item Brown
\end{enumerate}

\begin{center}
%\includegraphics[width=0.6\textwidth]{https://www.tech-faq.com/wp-content/uploads/images/eia-tia-568a-568b.png}
\end{center}

\section{RS-232 Cable}
The \textbf{RS-232 Cable} is used for serial communication between devices. It was commonly used for connecting PCs to modems and is still used in networking hardware for console connections.

\begin{center}
%\includegraphics[width=0.5\textwidth]{https://upload.wikimedia.org/wikipedia/commons/0/07/Serial_Cable.jpg}
\end{center}

\section{Fiber Optic (Duplex) Cable}
A \textbf{Fiber Optic Duplex Cable} consists of two fiber optic strands in a single cable, used for bidirectional data transfer. It allows high-speed and long-distance data communication.

\begin{center}
%\includegraphics[width=0.6\textwidth]{https://www.fs.com/images/fs/images2019/product/duplex-fiber-optic-cable.jpg}
\end{center}

\section{SFP Module}
The \textbf{Small Form-factor Pluggable (SFP)} module is a hot-swappable transceiver used in networking devices for fiber optic and Ethernet connections.

\begin{center}
%\includegraphics[width=0.5\textwidth]{https://media.fs.com/images/community/uploads/20200424/1-sfp-transceiver.jpg}
\end{center}

\section{Switch}
A \textbf{Switch} is a network device that connects devices in a LAN and uses MAC addresses to forward data to the correct destination.

\begin{center}
%\includegraphics[width=0.5\textwidth]{https://media.cablematters.com/images/thumbnails/515/515/detailed/5/2020011106560335.jpg}
\end{center}

\section{Router Configuration}
A \textbf{Router} connects different networks and directs data between them. Configuration typically involves setting IP addresses, routing protocols, and NAT.

\textbf{Example CLI Commands:}
\begin{verbatim}
interface fastethernet0/0
ip address 192.168.1.1 255.255.255.0
no shutdown
exit
\end{verbatim}

\begin{center}
%\includegraphics[width=0.4\textwidth]{https://cdn.ttgtmedia.com/rms/misc/router-wifi-mobile.png}
\end{center}

\section{Switch to Router Connection}
\textbf{Switch to Router Connection} is usually done through a straight-through Ethernet cable. The router's LAN port connects to the switch, enabling device communication and internet access.

\begin{center}
%\includegraphics[width=0.5\textwidth]{https://i.stack.imgur.com/sGb5w.png}
\end{center}

\section{OSI Model}
The \textbf{OSI (Open Systems Interconnection) Model} is a conceptual framework that standardizes the functions of a telecommunication or computing system into seven abstraction layers:
\begin{enumerate}
\item Physical
\item Data Link
\item Network
\item Transport
\item Session
\item Presentation
\item Application
\end{enumerate}

\begin{center}
%\includegraphics[width=0.6\textwidth]{https://upload.wikimedia.org/wikipedia/commons/thumb/d/dd/Osi-model-jb.svg/1920px-Osi-model-jb.svg.png}
\end{center}

\end{document}
